% Generated draft: requires mathematical and source review.
\documentclass[11pt]{article}
\usepackage[margin=25mm]{geometry}
\usepackage{fontspec}
\setmainfont{Latin Modern Roman}
\usepackage{amsmath,amssymb,mathtools,unicode-math}
\setmathfont{Latin Modern Math}
\usepackage{enumitem,needspace,xurl,hyperref}
\setlength{\parindent}{0pt}
\setlength{\parskip}{5pt}
\setlength{\emergencystretch}{3em}
\newcommand{\N}{\mathbb{N}}
\newcommand{\Z}{\mathbb{Z}}
\newcommand{\Q}{\mathbb{Q}}
\newcommand{\R}{\mathbb{R}}
\newcommand{\C}{\mathbb{C}}
\newcommand{\F}{\mathbb{F}}
\newcommand{\A}{\mathbb{A}}
\newcommand{\PP}{\mathbb{P}}
\newcommand{\E}{\mathbb{E}}
\newcommand{\eps}{\varepsilon}
\newcommand{\dd}{\,\mathrm{d}}
\begin{document}
\section*{UnsolvedMath ID: 800118}
\textbf{Source problem code:} catalogue entry
\subsection{Definitions and mathematical statement}
A closed manifold is compact and has no boundary. For an oriented circle bundle \(p\colon E\to M\), its integral Euler class \(e(p)\in H^2(M;\Z)\) is the obstruction to a section. Its square is the cup product
\[
e(p)^2=e(p)\smile e(p)\in H^4(M;\Z).
\]
A flat circle bundle has local trivializations whose transition maps are locally constant as functions of the base, with values in the chosen group of circle automorphisms. The supplied record does not specify that group or its regularity category.

Under the interpretation that the source symbol \(\chi\) denotes the integral Euler class of an oriented circle bundle, the requested construction is an oriented closed \(4\)-manifold \(M\) and a flat oriented circle bundle \(p\colon E\to M\) such that
\[
\chi=e(p),\qquad \chi^2=e(p)\smile e(p)\ne 0\quad\text{in }H^4(M;\Z).
\]
This formulation is conditional: the supplied record does not define \(\chi\), explicitly require orientation of the circle fibres, or specify whether the transition maps are homeomorphisms, diffeomorphisms, or elements of another structure group. These choices require clarification from the source context. In particular, flatness is not here silently interpreted as flatness of a principal \(S^1\)-bundle. No differentiability hypothesis on \(M\) is added.
\subsection{Short English statement}
Construct a flat circle bundle over an oriented closed \(4\)-manifold for which the source class \(\chi\) has nonzero integral cup square. If \(\chi\) is the Euler class, this asks for \(e(p)^2\ne0\). The meaning of \(\chi\), fibre orientation, and the allowed structure group remain unresolved in the supplied record.
\subsection{Sources}
\begin{itemize}
\item Local UnsolvedMath registry, permanent ID \(800118\), exact source problem code \texttt{catalogue entry}, titled \emph{Foliations and Diffeomorphism Groups --- Question 59}. The registry attributes the problem to Nariman and Yazdi, with proposed year 2025, and labels it open; this status is supplied metadata, not an independently verified assessment.
\item Nariman and Yazdi, \emph{Problem List on Foliations and Diffeomorphism Groups} (2025), Question 59. Supplied source URL: \url{https://arxiv.org/abs/2503.18274}. The record reports access on 2026-10-16 and extraction from arXiv LaTeX. These provenance details are supplied and have not been independently verified here.
\end{itemize}
\end{document}
